% Chapter 3: Randomized Experiments
% A First Course in Causal Inference - Peng Ding
% Rewritten in Stein motivation-first style with textbook formulas

\chapter{完全随机实验与费舍尔随机化检验}

\section{从潜在结果到随机实验}

在上一章中，我们建立了潜在结果框架，学会了如何清晰地定义因果效应。平均因果效应定义为：
\[
\tau = \mathbb{E}[Y_i(1) - Y_i(0)] = \bar{Y}(1) - \bar{Y}(0)
\]

但我们也清楚地看到了因果推断的根本问题：对于每个个体，我们只能观察到一种潜在结果，$Y_i(1)$ 和 $Y_i(0)$ 无法同时被观测。

那么，一个自然的问题是：我们能否设计某种数据收集过程，使得即使无法观察反事实，我们仍然可以估计因果效应？

这个问题的答案是肯定的，而解决方案就是\textbf{随机实验（Randomized Experiment）}——统计学的黄金标准。

\section{Fisher 的洞见：随机化的两个原则}

Ronald Fisher 在 1935 年的开创性著作《Design of Experiments》中指出了随机化的两个优势：

\begin{enumerate}
  \item \textbf{可比性}：它创造了在平均意义上可比的 treatment 和 control 组。
  \item \textbf{推理基础}：它作为统计推断的"合理基础"。
\end{enumerate}

第一点直观易懂——随机治疗分配不会偏向 treatment 或 control。第二点更微妙：Fisher 的意思是，随机化为统计检验提供了正当理由，这就是现在所谓的\textbf{费舍尔随机化检验（Fisher Randomization Test, FRT）}。

\section{完全随机实验的定义}

考虑一个包含 $n$ 个单元的实验，其中 $n_1$ 个接受治疗，$n_0$ 个接受对照（$n = n_1 + n_0$）。

\begin{Definition}[CRE]\label{def:cre}
固定 $n_1$ 和 $n_0$（其中 $n = n_1 + n_0$）。完全随机实验的治疗 assignment 机制为：
\[
\mathbb{P}(\bm{Z} = \bm{z}) = 1 \Big / \binom{n}{n_1},
\]
其中 $\bm{z} = (z_1,\ldots, z_n)$ 满足 $\sumn z_i = n_1$ 和 $\sumn (1-z_i) = n_0$。
\end{Definition}

在这个定义中，我们假设潜在结果向量 $\bm{Y}(1) = (Y_1(1), \ldots, Y_n(1))$ 和 $\bm{Y}(0) = (Y_1(0), \ldots, Y_n(0))$ 都是固定的。即使我们将它们视为随机的，我们可以在其上条件化，治疗分配机制变为：
\[
\mathbb{P}\{\bm{Z} = \bm{z} \mid \bm{Y}(1), \bm{Y}(0)\} = 1 \Big / \binom{n}{n_1}
\]
因为在 CRE 中 $\bm{Z} \ind \{\bm{Y}(1), \bm{Y}(0)\}$。

\textbf{关键含义}：在 CRE 中，$\bm{Z}$ 是从 $n_1$ 个 1 和 $n_0$ 个 0 的随机排列中产生的。

\section{Sharp Null 假设与 FRT 逻辑}

\citet{Fisher:1935} 感兴趣检验的零假设是：
\[
\HF: Y_i(1) = Y_i(0) \text{ 对所有单元 } i=1,\ldots, n.
\]
\citet{Rubin:1980} 称其为\textbf{sharp null 假设}，因为它可以基于观测数据确定所有潜在结果：$\bm{Y}(1) = \bm{Y}(0) = \bm{Y} = (Y_1, \ldots, Y_n)$，即观测结果向量。它也被称为\textbf{强零假设}。

概念上，在 $\HF$ 下，FRT 适用于任何检验统计量：
\begin{equation}\label{eq:teststatistic}
T = T(\bm{Z}, \bm{Y}) ,
\end{equation}
它是观测数据的函数。观测结果向量 $\bm{Y}$ 在 $\HF$ 下是固定的，所以 $T$ 中唯一的随机成分是治疗向量 $\bm{Z}$。实验者决定 $\bm{Z}$ 的分布，进而决定 $T$ 的分布——这就是计算 p 值的基础。

在 CRE 中，$\bm{Z}$ 在集合 $\{\bm{z}^1, \ldots, \bm{z}^M\}$ 上是均匀分布的，其中 $M = \binom{n}{n_1}$，$\bm{z}^m$ 是所有可能的包含 $n_1$ 个 1 和 $n_0$ 个 0 的向量。

\section{随机化分布与精确 p 值}

如果较大的 $T$ 值更极端，我们可以使用以下尾概率来衡量检验统计量相对于其随机化分布的极端性：
\begin{equation}\label{eq:exact-pvalue}
p_\frt = M^{-1} \sum_{m=1}^M \mathbb{I}\{T(\bm{z}^m,\bm{Y}) \geq T(\bm{Z},\bm{Y}) \}.
\end{equation}

这就是 Fisher 的 p 值。重要的是，$p_\frt$ 在 \cref{eq:exact-pvalue} 中对任何检验统计量和任何结果生成过程都有效。它在 $\HF$ 下是有限样本精确的（finite-sample exact）：
\begin{equation}\label{eq:p-value-valid}
\mathbb{P}(p_\frt \leq u) \leq u \quad \text{for all} \quad 0 \leq u \leq 1.
\end{equation}

\userannotation{为什么叫"有限样本精确"？}{这意味着 p 值的分布（虽然在离散情况下是阶梯函数）满足上式，因此 p 值在 $[0,1]$ 上均匀分布（对于连续情况）。}

\section{Monte Carlo 近似}

在实践中，$M$ 通常太大（例如当 $n=100, n_1=50$ 时，$M > 10^{29}$），枚举所有可能的治疗向量在计算上是不可行的。我们经常用 Monte Carlo 来近似 $p_\frt$：
\begin{equation}\label{eq:mc-pvalue}
\hat p_\frt  =  R^{-1} \sum_{r=1}^R   \mathbb{I}\{T(\bm{z}^r,\bm{Y})   \geq T(\bm{Z},\bm{Y})   \},
\end{equation}
其中 $\bm{z}^r$ 是 $\bm{Z}$ 的 $R$ 次随机置换。由于计算 $p_\frt$ 涉及对 $\bm{Z}$ 的置换，FRT 有时在 CRE 语境下被称为\textbf{置换检验（permutation test）}。

\textbf{修正的有限样本有效近似}：
\[
\tilde p_\frt  =  \frac{  1 + \sum_{r=1}^R   \mathbb{I}\{T(\bm{z}^r,\bm{Y})   \geq T(\bm{Z},\bm{Y})   \}  }{1+R}
\]
这个修正在任意有限的 $R$ 下都是有限样本有效的 p 值。

\section{检验统计量的规范选择}

FRT 的一个特点是它对任何检验统计量都生成有限样本精确 p 值。但我们应该选择能够给出 $\HF$ 可能违背信息的检验统计量。

\subsection{差值均值统计量}

差值均值统计量是：
\[
\hat{\tau}   = \hat{\bar{Y}}(1) - \hat{\bar{Y}}(0)
\]
其中
\[
\hat{\bar{Y}}(1)  =  n_1^{-1} \sum_{Z_i=1} Y_i  = n_1^{-1} \sumn Z_iY_i
\]
是 treatment 组的样本均值，
\[
\hat{\bar{Y}}(0) = n_0^{-1}  \sum_{Z_i=0} Y_i   = n_0^{-1}  \sumn (1-Z_i) Y_i
\]
是 control 组的样本均值。

在 $\HF$ 下，它有均值 $0$ 和方差：
\[
\text{var}(\hat{\tau})  = \frac{n}{n_1 n_0} s^2,
\]
其中
\[
s^2 = (n-1)^{-1} \sumn (Y_i - \bar{Y})^2, \quad \bar{Y} = n^{-1} \sumn Y_i.
\]

\textbf{渐近分布}：由于有限总体中心极限定理，随机化分布近似正态：
\begin{equation}\label{eq:frt-tau-hat}
\frac{ \hat{\tau} }{ \sqrt{\frac{n}{n_1 n_0} s^2} } \rightarrow \mathcal{N}(0,1)
\end{equation}
in distribution.

\textbf{与经典 t 检验的联系}：在 $\HF$ 下，近似 p 值接近经典两样本 t 检验（假设等方差）的 p 值。

基于代数运算（见习题 \ref{ex:3.8}），我们有如下展开式：
\begin{equation}\label{eq:frt-algebra}
(n-1)s^2 = \sum_{Z_i=1} \{Y_i - \hat{\bar{Y}}(1)\}^2 + \sum_{Z_i=0} \{Y_i - \hat{\bar{Y}}(0)\}^2 + \frac{n_1 n_0}{n} \hat{\tau}^2 .
\end{equation}

\userannotation{为什么 FRT 能证明 t 检验的合理性？}{传统 t 检验需要正态性假设。但 FRT 框架表明，只要随机化分配正确执行，t 检验程序在有限总体 + 随机化情境下可以使用——渐近理论只依赖于有限总体 CLT，而不依赖于数据的原始分布。}

\subsection{学生化统计量}

差值均值隐含地假设两组方差相等。如果治疗可能影响结果的变异性（异方差），我们需要使用学生化统计量：
\[
t = \frac{ \hat{\bar{Y}}(1) - \hat{\bar{Y}}(0) }{ \sqrt{ \frac{ \hat{S}^2(1) }{n_1} + \frac{\hat{S}^2(0)}{n_0} } },
\]
其中
\[
\hat{S}^2(1) = (n_1 - 1)^{-1} \sum_{Z_i =1} \{ Y_i -\hat{\bar{Y}}(1) \}^2,\quad
\hat{S}^2(0) = (n_0 - 1)^{-1} \sum_{Z_i =0} \{ Y_i -\hat{\bar{Y}}(0) \}^2
\]
分别是 treatment 组和 control 组的样本方差。

在 $\HF$ 下，有限总体 CLT 再次意味着 $t$ 渐近正态：$t \rightarrow \mathcal{N}(0,1)$。

这个统计量与 \texttt{t.test(var.equal = FALSE)} 的 p 值接近。

\userannotation{学生化为什么更稳健？}{当 treatment 组和 control 组的方差差异较大时，差值均值的检验可能失效。学生化通过分别估计两组的方差，能够更好地处理异方差情况。}

\subsection{Wilcoxon 秩和统计量}

差值均值和学生化统计量都关注均值差异，且对异常值敏感。Wilcoxon 秩和统计量基于合并观测结果的秩：
\[
W = \sumn Z_i R_i,
\]
其中 $R_i$ 是 $Y_i$ 在合并样本中的秩。

在 $\HF$ 下，$W$ 有均值 $\frac{n_1(n+1)}{2}$ 和方差 $\frac{n_1 n_0 (n+1)}{12}$，且渐近正态。

这个统计量的优点是：
\begin{itemize}
  \item 对异常值稳健
  \item 只依赖于数据的秩次信息
  \item 在许多备择假设下都有较好的功效
\end{itemize}

\subsection{Kolmogorov-Smirnov 统计量}

上述统计量主要关注均值差异。KS 统计量衡量两个分布之间的最大距离：
\[
D = \max_y |\hat{F}_1(y) - \hat{F}_0(y)|,
\]
其中 $\hat{F}_1(y) = n_1^{-1} \sumn Z_i I(Y_i \leq y)$。

这个统计量能够检测分布形状的差异，而不仅仅是均值差异。

\section{FRT 与传统 t 检验的关系}

一个重要的理论发现是：

\begin{center}
\textbf{FRT 为传统 t 检验提供了更坚实的基础。}
\end{center}

传统 t 检验假设"数据来自正态分布的无限总体"。但 FRT 框架表明，只要随机化分配正确执行，t 检验程序可以在"有限总体 + 随机化"的情境下使用。

这意味着，在随机实验中，我们可以不那么担心正态性假设——CLT 会为我们处理这个问题。

\section{LaLonde 实验数据案例}\label{sec:lalonde-data}

1986 年，Robert LaLonde 发表了一项影响深远的研究，探讨职业培训项目（National Supported Work Demonstration）对参与者收入的影响。

使用 R 进行 FRT 分析：

\begin{verbatim}
library(Matching)
data(lalonde)
z <- lalonde$treat
y <- lalonde$re78
\end{verbatim}

观测到的检验统计量值：

\begin{center}
\begin{tabular}{c|c|c}
\hline
统计量 & 观测值 & 描述 \\
\hline
$\hat{\tau}$ & 2.84 & 差值均值（t 统计量）\\
$t$ & 2.67 & 学生化统计量 \\
$W$ & 27402.5 & Wilcoxon 秩和 \\
$D$ & 0.132 & KS 统计量 \\
\hline
\end{tabular}
\end{center}

通过随机置换治疗向量，我们可以得到检验统计量的随机化分布的 Monte Carlo 近似。

Monte Carlo 置换检验的结果（$R = 10^4$）：

\begin{verbatim}
> exact.pv = c(mean(Tauhat >= tauhat),
+               mean(Student >= student),
+               mean(Wilcox >= W),
+               mean(Ks >= D))
> round(exact.pv, 3)
[1] 0.002 0.002 0.006 0.040
\end{verbatim}

所有单侧 p 值都小于 0.05——这提供了治疗效应显著的强有力证据。

\section{历史侧边栏：随机实验的起源}\label{sec:tea-history}

FRT 的逻辑可以追溯到早期的历史实例。

\subsection{James Lind 的坏血病实验}

James Lind (1716-1794) 是苏格兰医生和皇家海军卫生学的先驱。他进行了最早有清晰文献记录的随机实验之一，研究柑橘类水果对坏血病的影响。在 12 名坏血病患者被分配到 6 组后，接受柑橘类水果的组在六天后康复，而其他组没有。

如果将治疗简化为二元指示变量，观测到的极端 $2 \times 2$ 表格在 $\HF$ 下出现的概率为：
\[
p_\frt = \frac{1}{\binom{12}{2}} = \frac{1}{66} = 0.015.
\]

\subsection{Fisher 的女士品茶实验}

\citet{Fisher:1935} 描述了一个著名实验，一位女士声称能区分先加茶还是先加牛奶制成的奶茶。他制作了 8 杯茶（4 杯先加茶，4 杯先加牛奶），随机顺序呈现给女士。在她的全部猜对（$X = 4$）的情况下，p 值为：
\[
p_\frt = \frac{1}{70} = 0.014.
\]

这些例子强调了 Fisher 的两个原则：
\begin{enumerate}
  \item \textbf{随机化}：实验的随机分配保证了 p 值的有效性
  \item \textbf{重复}：必须有足够的单元以确保 FRT 有功效
\end{enumerate}

\section{Sharp Null 假设的置信区间}

FRT 的逻辑不仅限于 $\HF$。我们还可以检验更一般的 sharp null 假设：
\[
H_0(\bm{\tau}): Y_i(1) - Y_i(0) = \tau_i \quad \text{对所有 } i.
\]

通过检验与置信集的对偶性，我们可以构建 $\tau$ 的 $(1-\alpha)$ 置信集。

\section{本章小结}

本章建立了随机实验的理论与实践：

\begin{enumerate}
  \item \textbf{CRE 定义}：$\mathbb{P}(\bm{Z} = \bm{z}) = 1 / \binom{n}{n_1}$
  \item \textbf{随机化的两个原则}：可比性 + 推理基础
  \item \textbf{Sharp null}：$\HF: Y_i(1) = Y_i(0)$ 对所有 $i$
  \item \textbf{精确 p 值}：$p_\frt = M^{-1} \sum_{m=1}^M \mathbb{I}\{T(\bm{z}^m, \bm{Y}) \geq T(\bm{Z}, \bm{Y})\}$
  \item \textbf{多种统计量}：差值均值 $\hat{\tau}$、学生化 $t$、Wilcoxon $W$、KS $D$
  \item \textbf{理论联系}：FRT 为传统 t 检验提供基础
  \item \textbf{历史起源}：Lind 的坏血病实验、Fisher 的女士品茶
\end{enumerate}

下一章我们将学习 Neymanian 推断——另一种基于渐近理论的方法，它可以更方便地构建置信区间。

\section{习题}

\begin{Exercise}{3.1 — $p_\frt$ 的精确性}
\label{ex:3.1}
证明在 sharp null 下，$p_\frt$ 满足：
\[
\mathbb{P}(p_\frt \leq u) \leq u \quad \text{for all} \quad 0 \leq u \leq 1.
\]
\end{Exercise}

\begin{Exercise}{3.2 — Monte Carlo 误差}
\label{ex:3.2}
设 $p_\frt$ 是固定值，$\hat p_\frt$ 是其 Monte Carlo 估计（见 \cref{eq:mc-pvalue}）。证明：
\[
\mathbb{E}_{\text{mc}} (\hat p_\frt) = p_\frt, \quad
\text{var}_{\text{mc}} (\hat p_\frt) \leq \frac{1}{4R},
\]
其中下标 "mc" 表示由于 Monte Carlo 的随机性。
\end{Exercise}

\begin{Exercise}{3.3 — 有限样本有效的 Monte Carlo 近似}
\label{ex:3.3}
证明修正的 Monte Carlo 近似
\[
\tilde p_\frt  =  \frac{ 1 + \sum_{r=1}^R \mathbb{I}\{T(\bm{z}^r,\bm{Y}) \geq T(\bm{Z},\bm{Y}) \} }{1+R}
\]
在任意有限的 $R$ 下都是有限样本有效的 p 值。

\textbf{提示}：可以利用以下两个基本概率结果：
\begin{enumerate}
  \item 对于两个二项分布 $X_1 \sim \text{Binomial}(R, p_1)$ 和 $X_2 \sim \text{Binomial}(R, p_2)$，若 $p_1 \geq p_2$，则 $\mathbb{P}(X_1 \leq x) \leq \mathbb{P}(X_2 \leq x)$。
  \item 若 $p \sim \text{Uniform}(0,1)$ 且 $X \mid p \sim \text{Binomial}(R, p)$，则 $X$ 在 $\{0, 1,\ldots, R\}$ 上均匀分布。
\end{enumerate}
\end{Exercise}

\begin{Exercise}{3.4 — Fisher 精确检验}
\label{ex:3.4}
考虑二元结局的 CRE，数据汇总为 $2 \times 2$ 列联表：

\begin{center}
\begin{tabular}{lccc}
\hline
        & $Y =1$ & $Y =0$& 合计\\\hline
$Z =1$ & $n_{11}$ & $n_{10}$ & $n_1$\\
$Z =0$ & $n_{01}$ & $n_{00}$ & $n_0$\\
\hline
\end{tabular}
\end{center}

在 $\HF$ 下，证明任何检验统计量 $T(n_{11}, n_{10}, n_{01}, n_{00})$ 是 $n_{11}$ 的函数，且 $n_{11}$ 的精确分布是超几何分布。请说明超几何分布的参数。
\end{Exercise}

\begin{Exercise}{3.5 — 女士品茶详细计算}
\label{ex:3.5}
回忆女士品茶实验（见 \cref{sec:tea-history}）。计算 $\mathbb{P}(X = k)$，其中 $X$ 是女士正确识别 $k$ 杯先加牛奶的茶的个数，$k = 0, 1, 2, 3, 4$。
\end{Exercise}

\begin{Exercise}{3.6 — 协变量调整的 FRT}
\label{ex:3.6}
使用 LaLonde 数据集（见 \cref{sec:lalonde-data}），添加协变量进行 FRT 分析。假设所有潜在结果和协变量都是固定数值。

有两种策略：
\begin{enumerate}
  \item 运行结局对协变量的线性回归，得到残差（即"伪结局"）。基于残差定义四种检验统计量，进行 FRT 并报告 p 值。
  \item 定义检验统计量为结局对治疗和协变量线性回归中治疗系数的 FRT。
\end{enumerate}

解释为什么上述两种策略的五个 p 值都是有限样本精确的。
\end{Exercise}

\begin{Exercise}{3.7 — FRT 与广义线性模型}
\label{ex:3.7}
使用与 \cref{ex:3.6} 相同的数据集，但将结局改为 \texttt{re78 > 0} 的二元指示变量。运行结局对治疗和协变量的 logistic 回归。治疗系数是否显著？FRT p 值是多少？
\end{Exercise}

\begin{Exercise}{3.8 — 代数细节}
\label{ex:3.8}
验证 \cref{eq:frt-algebra}。
\end{Exercise}

\begin{Exercise}{3.9 — 推荐阅读}
\label{ex:3.9}
\citet{bind2020possible} 是一篇倡导在复杂实验中同时报告 p 值及其相应随机化分布的论文。
\end{Exercise}

% === 用户问答记录 ===
% (后续通过 interactive-qa 更新)
